% Synced from content.md by sync_latex.py.
\ifteachernotes
\section{\texorpdfstring{辅助线总结}{辅助线总结}}

\begin{center}
\small
\begin{tabularx}{\linewidth}{LLL}
\toprule
看见的条件或目标 & 可以先作的线 & 希望得到什么 \\
\midrule
已知外角平分线 & 从 $A$ 向 $BD$、$CD$ 作双垂 & 到角两边的距离相等 \\
外角相等，想造对称全等 & 延长 $CD$，取 $DE=DB$；或延长 $BD$，取 $DF=DC$ & $\mathrm{SAS}$ 全等，再借等腰换角 \\
单个 $45^\circ$ 或 $135^\circ$ & 过 $A$ 作 $AD$ 的垂线，交相应边或延长线 & 等腰直角三角形，接手拉手 \\
等腰直角 $ABC$，需搬运长度 & 从 $B$、$C$ 向 $AD$ 作双垂 & $K$ 型全等，直角边交叉对应 \\
目标为二倍关系 & 倍长短段，或给长段找中点 & 将二倍关系转为整段相等或半段相等 \\
例 $2$ 的角平分线 & 平行线补等腰；垂线取中点 & 把角平分线条件转成中点与等长 \\
\bottomrule
\end{tabularx}
\end{center}

两种“双垂”要分清：\textbf{$A$ 向 $BD$、$CD$ 作双垂，是利用角平分线的等距离；$B$、$C$ 向 $AD$ 作双垂，是利用等腰直角 $ABC$ 的 $K$ 型全等。} 它们的出发点和全等对应顺序都不同。

\needspace{4\baselineskip}\subsection{\texorpdfstring{使用时注意}{使用时注意}}

\begin{enumerate}
\item $DA$ 平分的是 $\angle BDX$，参与比较的两个角是 $\angle BDA$ 与 $\angle ADX$。
\item 条件与结论交换后，重新检查全等判定。等边待证时不能用它作 $\mathrm{HL}$ 条件；直角待证时不能提前用手拉手。
\item $K$ 型全等推出的是 $BM=AN$、$AM=CN$。写完对应边后，再按点的顺序加减公共部分，不能把同一条垂线上的两段凭图认成相等。
\end{enumerate}
\fi
